% Physical substrate extensions: F_21, lifting theorems, QCD predictions
% Physical substrate: F_21 QCD structure, lifting theorems, hadron spectroscopy

\section{Physical Substrate: \texorpdfstring{$F_{21}$}{F21}, Lifting, and QCD Structure}
\label{sec:physical_substrate}

Here $F_{21} = \mathbb{Z}_7 \rtimes \mathbb{Z}_3$ is the Frobenius group of order~21 (established in P35~\cite{SpivackGTEUnification}), and $\Phi_{\rm MDL}$ is the $\mathbb{Z}_7$-symmetric Klein--Gordon gauge field that is the continuum limit of $f_{\rm MDL}$ (P34~\cite{SpivackGTEMobius}, P35~\cite{SpivackGTEUnification}).
The discrete $f_{\rm MDL}$ Rule~110 substrate and its continuous $\Phi_{\rm MDL}$ field-theory limit share a common internal symmetry. The identification sharpens from a product $\mathbb{Z}_3 \times \mathbb{Z}_7$ to the finite non-abelian group
\begin{equation}
F_{21} = \mathbb{Z}_7 \rtimes \mathbb{Z}_3 \subset \mathrm{SU}(3),
\label{eq:f21_def}
\end{equation}
the Frobenius group of order~21 (denoted $\Sigma(21)$ in the finite-subgroup literature). The substrate field $\varphi$ remains a real scalar; what changes is the algebraic identity of the discrete internal symmetry. The abelianization $F_{21}^{\mathrm{ab}} = \mathbb{Z}_3$ is exactly the existing colour group $\{1,2,4\} \subset \mathbb{Z}_7^*$, so all prior anchor results on colour charge are preserved.

\subsection{\texorpdfstring{$F_{21}$}{F21} substrate identification}
\label{sec:f21_substrate}

Machine certification establishes that $F_{21}$ is MDL-forced rather than postulated.
The adjoint of $\mathrm{SU}(3)$ decomposes under $F_{21}$ as
\begin{equation}
\mathbf{8} = \mathbf{1}' \oplus \mathbf{1}'' \oplus \mathbf{3} \oplus \bar{\mathbf{3}},
\label{eq:f21_branching}
\end{equation}
verified at machine precision from $F_{21}$ three-dimensional irrep averaging on the Gell-Mann basis. This branching maps to the existing six-gluon-mediated vertex catalog plus two Cartan fields.
Including the LEP three-gluon constraint saves at least 38~bits of description length over a $\mathbb{Z}_7 \times \mathbb{Z}_3$ product with an external $\mathrm{SU}(3)$ postulate.
Lean: \lean{frobenius\_f21\_order}, \lean{frobenius\_su3\_branching\_8}, \lean{FrobeniusCatALAnchorInvariance}
(\lean{SylowIndexCouplingHierarchy.lean}, zero \texttt{sorry}).

The colour sector is identified with the Sylow-3 subgroup $\{1,2,4\} \subset \mathbb{Z}_7^*$
(\lean{color\_subgroup\_is\_sylow3}, \lean{z7\_color\_subgroup\_closed}).
This identification holds at the effective-field-theory level for momenta below
$\Lambda_{\mathrm{GTE}} \approx 2.0\,\mathrm{GeV}$ (the seven-kink full-winding threshold in physical units; \S\ref{sec:eft_domain}).
The proposal that the GTE gauge field $A_\mu$ corresponds to the Cartan generator $\lambda_3$ of $\mathrm{SU}(3)$ Yang--Mills, and that $F_{21}$ flows to $\mathrm{SU}(3)$ Yang--Mills above $\Lambda_{\mathrm{GTE}}$, constitutes a UV-completion \emph{conjecture}; verification remains open work.

\subsection{The two levels are one algebraic structure}
\label{sec:two_levels}

The UGP Physics programme operates at two resolutions: the discrete $f_{\rm MDL}$ Rule~110 automaton at finite lattice resolution~$M$, and the continuous $\Phi_{\rm MDL}$ $\mathbb{Z}_7$-Klein--Gordon gauge field in the continuum limit $M \to \infty$.
Three Lean-certified theorems connect the levels.

\paragraph{Algebraic Lifting (beable $\to$ physical).}
A PSC-admissible beable $B : \mathrm{Fin}\,5 \to \mathrm{Fin}\,7$ with positive $[D]$-weight determines a physical observable at the Compton scale carrying the same $\mathbb{Z}_7$ winding, $\mathbb{Z}_3$ colour, and generation index as~$B$.
Charge quantisation, colour structure, the three-generation hierarchy, and S-matrix quantum numbers proved at the beable level therefore hold at the physically observable scale.
Without this theorem, the certified algebraic structure at the substrate cell scale would lack a formal connection to electrons and quarks; with it, every beable-level result is licensed as a physical prediction.

\begin{theorem}[Algebraic Lifting, $\CatAL$]
\label{thm:algebraic_lifting}
Let $B : \mathrm{Fin}\,5 \to \mathrm{Fin}\,7$ be PSC-admissible with $\mathrm{DWeight}\,B > 0$.
Then $B$ determines a physical observable at the Compton scale with the same $\mathbb{Z}_7$ winding, $\mathbb{Z}_3$ colour, and generation index as~$B$.
\begin{proof}
Machine-certified in \lean{LiftingTheorem.lean}: \lean{algebraic\_lifting\_theorem}, zero \texttt{sorry}.
\end{proof}
\end{theorem}

\noindent\textbf{Physical meaning.}
The Algebraic Lifting Theorem asserts existence, not construction.
Any finite collection of SM particle types consistent with the PSC winding constraint must appear as a CA configuration in $f_{\rm MDL}$, even when the corresponding cellular configuration is unreachable by any finite computational process from simple initial conditions.
The distinction is physically significant: using $\hbar c \approx 197\,\mathrm{MeV{\cdot}fm}$ and the GTE lattice calibration ($\mathrm{sim\_to\_fm} = 0.100\,\mathrm{fm/sim}$), the proton Compton wavelength $\lambda_C^p = \hbar c/(m_p c^2) \approx 197\,\mathrm{MeV{\cdot}fm}/938\,\mathrm{MeV} \approx 0.21\,\mathrm{fm}$ corresponds to roughly $2$ simulation cells at $M=7$, scaling to ${\sim}10^{19}$--$10^{20}$ cells per proton Compton wavelength under the fine-end working hypothesis $a \approx \ell_{\rm Planck}$ (a calibration reading of the certificate, not a GTE derivation; footnote~\ref{fn:planck_working_hypothesis})---placing direct CA construction of a single hadron at that reading far beyond any realizable computation.
At the MDL-saturated matter-sector reading $a = \hbar c/\Lambda_{\mathrm{GTE}} \approx 0.1\,\mathrm{fm}$ (\S\ref{sec:eft_domain})---the calibration used here---the same wavelength is a few cells, which is what makes the substrate-regulated lattice campaigns of P42~\cite{SpivackPhiMDLField} feasible.
The existence proof circumvents the fine-end construction barrier entirely: it proceeds from the winding algebra alone, without constructing any explicit CA configuration.
The SM gauge sector is algebraically latent within the CA substrate; it is not imposed from outside.

\paragraph{3D spatial lifting (composites).}
In the one-dimensional ring model, an exhaustive search over 140{,}000 candidate quark--antiquark rings returns zero PSC-admissible meson configurations: adjacent opposite-colour cells cross an orbit-family boundary forbidden by the PSC table.
In the three-dimensional $f_{\rm MDL}$ causal graph, a quark at node~$A$ and antiquark at node~$B$ connected by a vacuum path need never be adjacent; PSC-admissibility decomposes over individual nodes and vacuum cells (\lean{vacuum\_psc\_admissible}).

\begin{theorem}[3D spatial composite lifting, $\CatAL$]
\label{thm:spatial_lifting}
Let $B_A$, $B_B$ be PSC-admissible beables at distinct positions with
$\mathrm{color}(B_A) + \mathrm{color}(B_B) \equiv 0 \pmod{3}$ and an \lean{IsVacuumPath} between them.
Then the composite is PSC-admissible, colour-neutral, and a physical bound state.
\begin{proof}
\lean{spatially\_extended\_composite\_lifting}, \lean{meson\_bound\_state\_exists},
\lean{baryon\_bound\_state\_exists} in \lean{SpatiallyExtendedLifting.lean}, zero \texttt{sorry}.
Path existence \lean{causal\_path\_exists} is proved as a theorem for forward-causal pairs
($t_a \leq t_b$), with zero custom axioms in any \CatAL\ proof chain (\lean{SpatiallyExtendedLifting.lean}).
Meson and baryon existence are therefore unconditional for forward-causal quark--antiquark
configurations.
\end{proof}
\end{theorem}

The vacuum path between separated colour charges is the causal-graph structure corresponding to a gluon flux tube in QCD language.

\paragraph{Algebraic descent (continuous $\to$ discrete).}
Any algebraic or topological property of $\Phi_{\rm MDL}$ depending only on $F_{21} = \mathbb{Z}_7 \rtimes \mathbb{Z}_3$ holds for $f_{\rm MDL}$ at every resolution $M \geq 1$.
Asymptotic freedom ($b_0 = 7$), strong-CP resolution ($\theta_{\rm QCD} = 0$), Casimir invariants ($C_F = 4/3$, $C_A = 3$), colour confinement, and three-generation structure are \emph{algebraic} and therefore $M$-independent.
Geometric properties---Lorentz invariance, SR accuracy---require $M \to \infty$ and differ at finite~$M$ by the lattice correction $\varepsilon_0(M) = \pi^2/(3M^2)$ (at $M = 7$: $\varepsilon_0 \approx 6.71\%$).

\begin{theorem}[Algebraic Descent, $\CatAL$]
\label{thm:algebraic_descent}
Let $P$ be an algebraic or topological property of $\Phi_{\rm MDL}$ depending only on $F_{21}$.
Then $P$ holds for $f_{\rm MDL} = \Phi_{\rm MDL}|_{M,\,\mathrm{binary}}$ at every $M \geq 1$.
\begin{proof}
\lean{algebraic\_descent\_theorem} in \lean{AlgebraicDescentTheorem.lean}, zero \texttt{sorry}.
\end{proof}
\end{theorem}

\noindent\textbf{Physical meaning.}
The Algebraic Descent Theorem is the converse of the Lifting Theorem.
Where lifting maps discrete CA configurations upward to physical observables at the Compton scale, descent maps continuum properties downward to every discrete resolution.
Any conservation law, charge quantisation condition, or generation-structure result proved at the CA level in Lean---asymptotic freedom ($b_0 = 7$, $b_1 = 26$), colour confinement, $\theta_{\rm QCD} = 0$, exactly three generations---is not a discrete-model approximation but an exact statement about the continuum theory $\Phi_{\rm MDL}$, valid at every resolution $M \geq 1$.
The Lean certifications in Sections~\ref{sec:color_confinement}--\ref{sec:asymptotic_freedom} are therefore exact continuum constraints, not lattice artefacts.

\begin{remark}[Discrete--continuum algebraic correspondence, $\CatAL$]
\label{rem:discrete_continuum_correspondence}
Theorems~\ref{thm:algebraic_lifting} and~\ref{thm:algebraic_descent} together establish a two-directional algebraic correspondence between the discrete $f_{\rm MDL}$ CA and the continuous $\Phi_{\rm MDL}$ field theory in the $F_{21}$-winding sector.
\begin{itemize}
\item \emph{Upward (discrete $\to$ physical):} every PSC-admissible beable configuration lifts to a physical observable at the Compton scale with the same $\mathbb{Z}_7$ winding, colour, and generation index. The SM particle spectrum is algebraically contained in the CA substrate, independently of constructibility.
\item \emph{Downward (continuous $\to$ discrete):} every algebraic property of the continuum theory depending only on $F_{21}$ descends to every discrete resolution $M \geq 1$. CA-level conservation laws and charge assignments are exact continuum constraints.
\end{itemize}
The two directions together establish that the SM winding sector and the Rule~110/$f_{\rm MDL}$ CA substrate are algebraically isomorphic in this sector: they are not competing descriptions but two resolutions of one algebraic structure.
The SM is algebraically realised within Rule~110 in the sense that lifting places every PSC-admissible particle configuration there; Rule~110 is algebraically constrained by any continuum realisation of the same SM winding algebra in the sense that descent forces the CA conservation laws upon it.
The structure is summarised in Figure~\ref{fig:discrete_continuum_correspondence}.
Machine-certified: \lean{algebraic\_lifting\_theorem} and \lean{algebraic\_descent\_theorem}
(\lean{LiftingTheorem.lean}, \lean{AlgebraicDescentTheorem.lean}, zero \texttt{sorry}).
\end{remark}

\begin{figure}[h]
\centering
\resizebox{\textwidth}{!}{%
\begin{tikzpicture}[
  every node/.style={font=\small},
  box/.style={draw, rounded corners=4pt, minimum height=3.0cm, text width=2.6cm,
              align=center, inner sep=6pt},
  leftbox/.style={box, fill=blue!8, draw=blue!50!black, text width=3.1cm},
  centerbox/.style={box, fill=gray!10, draw=gray!55},
  rightbox/.style={box, fill=orange!8, draw=orange!55!black},
  arr/.style={-{Stealth[length=6pt,width=5pt]}, thick, shorten >=4pt, shorten <=4pt},
]
% Boxes — wider spacing (7 cm between centres) leaves ~3.7 cm for arrows
\node[leftbox]  (cont) at (0,0) {
  \textbf{Continuum SM}\\[3pt]
  $\mathrm{SU}(3){\times}\mathrm{SU}(2){\times}\mathrm{U}(1)$\\[1pt]
  gauge structure\\[4pt]
  particle content,\\charges, Lagrangian\\[4pt]
  conservation laws
};
\node[centerbox] (core) at (7.0,0) {
  \textbf{Algebraic Core}\\
  \small($\mathbb{Z}_7$ winding algebra)\\[3pt]
  SM generation orbit\\[3pt]
  PSC uniqueness\\constraint\\[3pt]
  $f_{\rm MDL}$ orbit structure
};
\node[rightbox]  (ca)   at (14.0,0) {
  \textbf{Discrete CA}\\
  \small(Rule~110, $f_{\rm MDL}$)\\[3pt]
  $\mathbb{Z}_7$ winding lattice\\[3pt]
  CUP-4 orbit forcing\\[3pt]
  CA conservation laws\\(Lean-certified)
};
% Descent arrow: algebraic core -> continuum
\draw[arr] (core.west) --
  node[midway, above, align=center, font=\small\bfseries] {Algebraic Descent}
  node[midway, below, align=center, font=\scriptsize\itshape] {CA theorems\\exact continuum $\Leftarrow$\\constraints}
  (cont.east);
% Lifting arrow: algebraic core -> discrete CA
\draw[arr] (core.east) --
  node[midway, above, align=center, font=\small\bfseries] {Algebraic Lifting}
  node[midway, below, align=center, font=\scriptsize\itshape] {SM configurations\\$\Rightarrow$ CA\\realisations}
  (ca.west);
% Overall correspondence label
\node[font=\small\itshape, below=0.85cm of core]
      {Discrete--Continuum Structural Correspondence};
\end{tikzpicture}%
}
\caption{The discrete--continuum structural correspondence. \emph{Algebraic Lifting} (rightward
arrow) proves that every PSC-admissible SM particle configuration is algebraically realised
in the CA substrate; \emph{Algebraic Descent} (leftward arrow) proves that every CA-level
conservation law holds in any continuum realisation of the same $F_{21}$ algebra.
Together they establish that the discrete and continuum descriptions of the SM winding sector
are algebraically isomorphic (Remark~\ref{rem:discrete_continuum_correspondence}).}
\label{fig:discrete_continuum_correspondence}
\end{figure}

\subsection{Physical consequences at the Compton scale}
\label{sec:physical_consequences}

The Lifting Theorem (\ref{thm:algebraic_lifting}) packages three structural SM predictions as machine-verified theorems rather than separate postulates.

\begin{itemize}
\item \textbf{Charge quantisation} in units of $1/3$
(\lean{charge\_quantization\_physical}, machine-certified in Lean~4).
\item \textbf{Exactly three generations} from the period-3 $\mathbb{Z}_7$ orbit
(\lean{three\_generations\_physical}, machine-certified in Lean~4).
\item \textbf{No fourth generation}: \lean{no\_fourth\_generation\_physical} exhausts all PSC-admissible orbit classes and shows no fourth-generation class exists; \lean{no\_exotic\_physical\_generation} rules out any observable generation outside $\{\mathrm{gen}_1,\mathrm{gen}_2,\mathrm{gen}_3\}$.
\item \textbf{S-matrix framework}: \lean{smatrix\_framework\_exists} bundles scattering existence, charge quantisation, S-matrix consistency, the absence theorem, and the three-generation proof as a coherent five-conjunct certificate.
\end{itemize}

The Absence Theorem (\lean{algebraic\_absence\_theorem}) rules out forbidden processes at the beable level; combined with \lean{scattering\_existence} (\lean{LiftingTheorem.lean}~\S11), every algebraically allowed SM scattering process from the vertex catalog exists at physical scale. Quantitative cross-sections require three-dimensional $f_{\rm MDL}$ Hamiltonian dynamics and remain beyond the present certification scope.

These results compose into the three-generation uniqueness capstone
(\lean{gte\_uniquely\_predicts\_three\_generations}, \lean{ThreeGenerationCapstone.lean},
$\CatAL$, zero \texttt{sorry}): every PSC-consistent substrate is record-equivalent to
$f_{\rm MDL}$, which admits exactly three distinct non-vacuum generation states and no
fourth generation at physical scale.

\subsection{Why \texorpdfstring{$\mathbb{Z}_7$}{Z7}?}
\label{sec:why_z7}

Two independent arguments force $\mathbb{Z}_7$ as the minimal arithmetic field.

\paragraph{Sylow-3 colour.}
$\mathbb{Z}_7^* \cong \mathbb{Z}_6$; its unique order-3 subgroup is $\{1,2,4\}$, simultaneously the Sylow-3 subgroup and the cubic roots of unity in GF(7) (\lean{color\_subgroup\_is\_sylow3}).

\paragraph{GF(7) minimality.}
$\mathbb{Z}_7$ is the smallest prime field containing both a subgroup of order~2 ($\mathbb{Z}_2$, parity) and a subgroup of order~3 ($\mathbb{Z}_3$, colour):
\lean{gf7\_minimal\_for\_z2\_z3} (zero \texttt{sorry}).
The three-generation orbit structure and colour triplet are forced by field minimality, not by independent postulates.

\begin{remark}[The $f_{\rm MDL}$ orbit is not a GF(7) group code]
\label{rem:not_gf7_group_code}
Although $\{1,2,4\} \subset \mathbb{Z}_7^*$ forms a group code (the unique order-3 subgroup, cubic roots of unity in $\mathrm{GF}(7)$), the full $f_{\rm MDL}$ generation orbit
$\{\mathrm{gen}_1, \mathrm{gen}_2, \mathrm{gen}_3, \mathrm{vacuum}\} \subset \mathbb{Z}_7^5$
is \emph{not} a GF(7) group code: additive closure fails modulo~7 (e.g.\ $\mathrm{gen}_1 + \mathrm{gen}_1 \notin \{\mathrm{gen}_1, \mathrm{gen}_2, \mathrm{gen}_3, \mathrm{vacuum}\}$
componentwise mod~7), and the orbit is not closed under the field operations.
The colour subgroup $\{1,2,4\} \subset \mathbb{Z}_7^*$ carries genuine GF(7) group structure; the full generation orbit does not.
Computationally verified across all pairwise sums and all standard code-structure tests (\texttt{gf7\_orbit\_code\_check.py}).
\end{remark}

\subsection{\texorpdfstring{$f_{\rm MDL}$}{fMDL} as a Wolfram Category~IV rule}
\label{sec:fmdl_cat4}

Three of four standard Category~IV diagnostic tests are positive for $f_{\rm MDL}$ with the actual lookup table (computationally verified).
Among $\mathbb{Z}_7$ cellular automata, $f_{\rm MDL}$ is the unique MDL-minimal Category~IV rule---the $\mathbb{Z}_7$ analog of Rule~110.
The mutual implication GTE $\leftrightarrow$ Rule~110 is the binary shadow (parity projection) of the deeper GTE $\leftrightarrow$ $f_{\rm MDL}$ relationship: observers mediated by $[D]$ see Rule~110; full $\mathbb{Z}_7$ observers see $f_{\rm MDL}$.

\subsection{Colour confinement at the beable level}
\label{sec:color_confinement}

\begin{theorem}[No free quarks, $\CatAL$]
\label{thm:no_free_quark}
No PSC-admissible single-quark beable exists in the GTE framework.
\begin{proof}
Exhaustive over $7^5 = 16{,}807$ states: \lean{no\_psc\_admissible\_single\_quark}
(\lean{ColorConfinement.lean}, zero \texttt{sorry}).
\end{proof}
\end{theorem}

All PSC-admissible composite systems are colour-neutral:
\lean{psc\_rc\_requires\_system\_color\_neutral} (machine-certified in Lean~4).
Isolated coloured generation states cascade to vacuum in exactly $N_{\rm gen} = 3$ steps in the $F_{21}$-extended CA (computationally verified, \texttt{z3\_z7\_color\_extension.py}).

At the QFT level, in the $\mathbb{Z}_3$-gauged $\Phi_{\rm MDL}$ theory, the lightest colour-singlet excitation has positive mass $\Delta_{\rm gauge} > 0$.
The analytical lower bound $\Delta_{\rm gauge} \geq 2 M_{\rm kink,lat} = 0.300$ in lattice units ($\approx 592\,\mathrm{MeV}$ at the Compton-scale calibration $\mathrm{sim\_to\_fm} = 0.100\,\mathrm{fm/sim}$) is certified unconditionally:
\lean{qft\_gauged\_mass\_gap\_unconditional}, \lean{qft\_gauged\_mass\_gap\_pos\_unconditional}
(\lean{QFT/GaugedMassGap.lean}, zero \texttt{sorry}).
The beable-level mass gap \lean{gte\_mass\_gap} ($\exists\,\Delta > 0$, machine-certified in Lean~4) and the QFT-level result compose in \lean{confined\_massive\_color\_singlet} and the capstone \lean{unified\_physical\_exclusion}: any $[D]$-weighted physical state is PSC-admissible, colour-confined, massive, and anomaly-free.
The mass gap has a physical lower bound: the lightest standard-model generation has mass
$\Delta \geq 1.8\,\mathrm{MeV}$ (conservative PDG lower bound on $m_u$), certified in Lean~4
(\lean{gte\_mass\_formula\_physical}, zero axioms, \lean{MassGap.lean}) \cite{SpivackEmergentGravity}.

\paragraph{Yang--Mills continuum programme.}
The objection that a $\mathbb{Z}_3$ product cannot support non-abelian $\mathrm{SU}(3)$ gauge structure is removed by $F_{21} \subset \mathrm{SU}(3)$: $F_{21}$ is a finite non-abelian subgroup of $\mathrm{SU}(3)$.
A rigorous mass-gap proof would proceed via (i)~pure-gauge $F_{21}$ mass gap (beable and QFT levels, machine-certified in Lean~4), (ii)~rigorous $F_{21} \to \mathrm{SU}(3)$ continuum limit (\emph{conjectural}), and (iii)~Osterwalder--Schrader / Wightman axiom verification (\emph{open}).
This is a multi-year programme; it does not block the EFT-level predictions below.

\subsection{Two distinct \texorpdfstring{$\mathbb{Z}_7$}{Z7} winding quantities}
\label{sec:wa_wb}

Two distinct $\mathbb{Z}_7$ quantities appear in the framework and must not be conflated.

\begin{itemize}
\item $W_A$ (\textbf{Level~A, orbit-sum beable charge}): the sum of $\mathbb{Z}_7$ labels in a five-cell orbit configuration.
\item $W_B$ (\textbf{Level~B, SM interaction winding}): the P22 CUP-4/11 winding label $v$ assigned to each particle for vertex conservation.
\end{itemize}

For gen$_1$ and gen$_2$ (sum-4 conserving class), $W_A = W_B$ coincidentally.
For gen$_3$, $W_A = 24 \equiv 3 \pmod{7}$ while $W_B(\tau) = 4$: this reflects orbit position versus interaction charge, not an internal inconsistency.
The $\tau$ lepton with $Q = (4,1)$ is confirmed at Level~B.

\subsection{Composite states: baryons and mesons}
\label{sec:composite_states}

Composite configurations with total winding $0 \bmod 7$ form the algebraic candidate space.

\paragraph{One-dimensional ring model.}
\textbf{Baryons confirmed:} 135 non-trivial PSC-admissible colour-neutral $k=3$ composites exist.
\textbf{Mesons null:} three independent methods---strict scan, gluon-bridge exhaustive search (0 true mesons in 140{,}000 rings), and de~Bruijn graph analysis---all return zero true mesons as static 1D rings.
A one-dimensional $f_{\rm MDL}$ collision test finds GEN1+GEN1 interactions are causal glider collisions, not vacuum-mediated binding; no bound state forms on an open tape.

\paragraph{Three-dimensional causal graph.}
Theorem~\ref{thm:spatial_lifting} resolves the apparent contradiction: mesons exist as spatially separated quark--antiquark pairs connected by vacuum paths, not as adjacent 1D rings.
Baryons exist both as 1D orbit composites (Rank~18 ring enumeration) and as 3D vacuum-path composites (\lean{PhysicalBaryonState}, \lean{baryon\_bound\_state\_exists}); the two results are complementary models of the same physical states.
In the $\mathbb{Z}_7$-KG continuum picture, mesons additionally exist as kink--antikink breathers with Yukawa (not linear) inter-kink force; the algebraic-only treatment captures the candidate space, while dynamics supplies the binding mechanism.

\subsection{Asymptotic freedom and the running coupling}
\label{sec:asymptotic_freedom}

From $F_{21}$ with $N_c = 3$ (three-dimensional irrep dimension) and $N_f = 6$ (GTE species formula $W_B = 4k \bmod 7$, $k \in \{4,5\}$ quarks, three colours each), the one-loop $\beta$-function coefficient is
\begin{equation}
b_0 = 7, \qquad \beta = -\frac{7 g^3}{16\pi^2},
\label{eq:beta_one_loop}
\end{equation}
matching QCD exactly.
Lean: \lean{f21\_substrate\_beta\_coefficient}, \lean{f21\_substrate\_asymptotic\_freedom} (machine-certified in Lean~4).
Leptons ($k=1$, colour singlets) do not contribute to the running; their $F_{21}$ representation is trivial.

At two loops,
\begin{equation}
b_1 = \frac{11}{3} N_c^2 - \frac{2}{3} N_f N_c - \frac{1}{6} N_f = 26,
\label{eq:beta_two_loop}
\end{equation}
again matching QCD; \lean{f21\_two\_loop\_beta\_coefficient} (machine-certified in Lean~4).
Predicted values: $\alpha_s(M_Z) = 0.1318$ at one loop (11.67\% above PDG) and $\alpha_s(M_Z) \approx 0.1201$ at two loops (+1.75\% vs PDG).
Full Dyson resummation of the running coupling remains \emph{conjectural}.
Both $\beta$ coefficients are conditional on the $F_{21} \subset \mathrm{SU}(3)$ substrate identification.

\subsection{Hadron spectroscopy from GTE kink composites}
\label{sec:hadron_spectroscopy}

Kink--antikink and multi-kink composites in the $\Phi_{\rm MDL}$ substrate reproduce hadron multiplet structure with zero PDG inputs in the derivation chain.
Confidence levels are indicated explicitly.

\begin{table}[ht]
\centering
\caption{Hadron spectroscopy predictions from GTE kink composites.
Values marked ``Lean-certified'' are zero sorry; ``comp.'' are computational; PROVISIONAL entries require continuum-limit confirmation.}
\label{tab:hadron_predictions}
\footnotesize
\begin{tabular}{@{}llcc@{}}
\toprule
Observable & GTE prediction & PDG / lattice & Status \\
\midrule
Meson nonet ($J^P=0^-$) & 9 states, correct $J^P$ & 9 PDG states & PROVISIONAL \\
Baryon octet ($J^P=\tfrac{1}{2}^+$) & 8 states & 8 PDG states & Lean-certified \\
Baryon decuplet ($J^P=\tfrac{3}{2}^+$) & 10 states & 10 PDG states & Lean-certified \\
Vector meson nonet ($J^P=1^-$) & RMS $\approx 3.2\%$ & $\rho,K^*,\omega,\phi$ & PROVISIONAL \\
Quark masses (6) & within 7\% & PDG current & computational \\
$r_s/d$ & 20.00 & 19.80 & computational \\
$f_\pi$ & 91.35 MeV & 92.07 MeV & computational \\
$m_\pi^\pm$ (GOR+condensate) & 139.57 MeV & 139.57 MeV & $\CatAL$ (\lean{pion\_mass\_from\_gor}; P39) \\
$\theta_P$ ($\eta$--$\eta'$ mixing) & $-13.08^\circ \pm 3.74^\circ$ & $[-14.3^\circ,-10.7^\circ]$ & computational \\
$\chi_{\rm top}^{1/4}$ & 166.5 MeV & 178 MeV (lattice) & PROVISIONAL \\
Strong CP & $\theta_{\rm QCD} = 0$ & consistent & Lean-certified \\
\bottomrule
\end{tabular}
\end{table}

The $\eta$--$\eta'$ mixing angle $\theta_P = -13.08^\circ \pm 3.74^\circ$ lies in the PDG range $[-14.3^\circ, -10.7^\circ]$.
The chain runs: GTE quark masses $\to$ $f_\pi = m_{\rm kink}/\pi$ (DHN/BPS, $-0.81\%$ vs PDG) $\to$ four-dimensional string tension $\sigma_{4\mathrm{D}}$ $\to$ topological susceptibility $\chi_{\rm top}^{1/4} = 166.5\,\mathrm{MeV}$ ($-6.4\%$ vs lattice QCD 178 MeV) $\to$ one-loop $B_0$ condensate $\to$ NLO quark self-energy correction $\to$ $\theta_P$.
Topological susceptibility uses the corrected formula with $\sigma_{4\mathrm{D}} = (N_3/N_7)\sigma_{2\mathrm{D}}$; an earlier 190.2 MeV estimate used incorrect density units and is superseded.

Second baryon octet states are suppressed by $\delta K = \log_2(3) = 1.585$ bits from QR(7) $\mathbb{Z}_3$ symmetry (\lean{gte\_physical\_baryon\_octet\_is\_ms}, machine-certified in Lean~4).
Spin-$1/2$ baryon assignment is MDL-forced (\lean{jpspin\_kinks\_are\_spin\_half}, machine-certified in Lean~4).
Strong CP $\theta_{\rm QCD} = 0$ follows from $F_{21}$ discrete holonomy (\lean{f21\_theta\_term\_vanishes}, machine-certified in Lean~4).

\subsection{Renormalizability, vertices, and decay rates}
\label{sec:renorm_vertices_decays}

\paragraph{Renormalizability is PSC-forced.}
Non-renormalizable operators (dimension $> 4$) require Zone~L2 beables (\lean{NonRenormalizable} $\equiv$ $\neg\,\mathrm{PSCAdmissible}$).
The SM Lagrangian is renormalizable as a PSC consequence:
\lean{renormalizability\_psc\_forced}, \lean{physical\_renormalizability}, \lean{sm\_lagrangian\_renormalizable}
(\lean{AnomalyRenormalizability.lean}, zero \texttt{sorry}).
Anomaly cancellation and per-generation charge neutrality bundle in \lean{all\_generation\_charge\_neutrality\_bundle}.

\paragraph{Vertex catalog.}
Under strict $\mathbb{Z}_3$ gauge invariance, all 7/7 GTE vertices are recovered with zero spurious vertices (three-tier proof: topological, algebraic, numerical; computationally verified, \texttt{scripts/gauged\_vertex\_catalog.py}).

\paragraph{Kink collision identification.}
Each entry in the 7/7 GTE vertex catalog corresponds to a topologically allowed kink-antikink scattering event in the $\Phi_{\rm MDL}$ field theory.
A kink-antikink pair in the $(\varphi, \chi)$ plane is uniquely labelled by the joint topological charge $(Q_\varphi, Q_\chi)$, where $Q_\varphi \in \mathbb{Z}_7$ is the $\mathbb{Z}_7$ winding number and $Q_\chi \in \mathbb{Z}_3$ is the colour charge; the four orbit states carry
$\mathrm{gen}_1 \leftrightarrow (4,1)$, $\mathrm{gen}_2 \leftrightarrow (4,2)$, $\mathrm{gen}_3 \leftrightarrow (3,1)$, $\mathrm{vacuum} \leftrightarrow (0,0)$.
A kink-antikink collision conserving both $Q_\varphi$ and $Q_\chi$ maps precisely to an entry in the GTE vertex catalog; enforcing joint conservation strictly refines the set of allowed transitions relative to $\mathbb{Z}_7$-winding conservation alone, eliminating no physical SM vertex (computationally verified, \texttt{scripts/gauged\_vertex\_catalog.py}).
The charged-current analog in the leptoquark sector is derived in Paper~P33~\cite{SpivackDeeperConsequences}.

\paragraph{Tree-level decay rates.}
GTE masses reproduce PDG to 0.4\%; tree-level partial widths give $\Gamma_\mu$ at $+12.2\%$ and $\Gamma_\tau$ at $+12.5\%$ relative to PDG, with $\tau_\tau/\tau_\mu$ ratio error 0.3\%.
Rate positivity and ordering are machine-certified via the Lifting Theorem (computationally verified; companion script \texttt{decay\_rates\_from\_gte.py} in Paper~33).
Radiative corrections are open.

\paragraph{GoE stability at physical scale.}
The Garden-of-Eden stability result lifts to the Compton scale via the Lifting Theorem, certifying first-generation stability as a physical observable property.

\subsection{Domain of validity and the GTE effective field theory}
\label{sec:eft_domain}

The GTE framework is an effective field theory below the matching scale
\begin{equation}
\Lambda_{\mathrm{GTE}} = 7 \cdot m_{\rm kink}^{\rm phys} \approx 2.0\,\mathrm{GeV},
\label{eq:lambda_gte}
\end{equation}
the seven-kink full-winding threshold: the mass of the minimal $\mathbb{Z}_7$-neutral chain of like-charge kinks of the non-compact substrate field, with multiplier $7 = |\mathbb{Z}_7|$ (the threshold is exact because like-charge kinks repel).
The tree (classical BPS) kink mass $m_{\rm kink} = (8/49)\,m_\tau$ gives $\Lambda_{\mathrm{GTE}} = (8/7)\,m_\tau = 2030.70 \pm 0.14\,\mathrm{MeV}$ ($\CatAD$); the quantum (pole) kink mass $M^Q = 281 \pm 21\,\mathrm{MeV}$ gives $7 M^Q = 1.970 \pm 0.146\,\mathrm{GeV}$ ($\CatA$).
The two readings agree within $0.5\sigma$; the envelope is $1.96 \pm 0.15\,\mathrm{GeV}$ (band $< 10\%$).
The derivation, the multiplier mechanism, and the machine-certified arithmetic core are given in P39~\cite{SpivackGTEQCDStructure}.
Above $\Lambda_{\mathrm{GTE}}$, the $F_{21} \to \mathrm{SU}(3)$ Yang--Mills UV-completion conjecture applies; below it, GTE plays the role that chiral perturbation theory ($\chi$PT) plays for pion physics below $\Lambda_{\chi\mathrm{SB}}$.

\begin{table}[ht]
\centering
\caption{Prediction classification by energy scale relative to $\Lambda_{\mathrm{GTE}} \approx 2.0\,\mathrm{GeV}$.}
\label{tab:eft_classification}
\small
\begin{tabular}{@{}p{0.38\linewidth}p{0.22\linewidth}p{0.28\linewidth}@{}}
\toprule
Prediction class & Energy regime & Confidence \\
\midrule
Charge quantisation, three generations, colour confinement (beable) & All scales (algebraic) & Lean-certified \\
$\beta$-function coefficients $b_0=7$, $b_1=26$ & UV (conditional on $F_{21}$) & Lean-certified \\
Hadron multiplets, $\theta_P$, $\chi_{\rm top}$ & IR, $E < \Lambda_{\mathrm{GTE}}$ & computational / PROVISIONAL \\
$\alpha_s(M_Z)$ two-loop & Intermediate & Lean-certified/PROVISIONAL \\
Cartan gluon $A_\mu \leftrightarrow \lambda_3$ & UV, $E > \Lambda_{\mathrm{GTE}}$ & Conjectural \\
Wightman/OS axiom bridge & Continuum limit & Open \\
\bottomrule
\end{tabular}
\end{table}

The formal matching condition at $\Lambda_{\mathrm{GTE}}$ equates GTE kink-composite observables below the threshold with $\chi$PT and QCD sum-rule observables above it; the UV-completion conjecture asserts that integrating out the discrete $F_{21}$ substrate yields $\mathrm{SU}(3)$ Yang--Mills with the certified $\beta$-function coefficients.
Neither conjecture is required for the IR hadronic predictions in Table~\ref{tab:hadron_predictions}, which are evaluated entirely within the GTE EFT.
